\section{Related Work}
\label{sec:related}

\paragraph{VLA policies and their benchmarks.} RT-1~\citep{rt1} showed that a transformer trained on many real demonstrations can control a robot across tasks, and RT-2~\citep{rt2} fine-tuned a web-scale vision-language model to emit actions as text tokens. Octo~\citep{octo} and OpenVLA~\citep{openvla} are open generalist policies trained on pooled robot datasets; OpenVLA builds on the Prismatic vision-language model~\citep{prismatic}, with DINOv2~\citep{dinov2} and SigLIP~\citep{siglip} image features and a Llama~2 backbone~\citep{llama2}. OpenVLA-OFT~\citep{oft} changed its fine-tuning recipe to parallel decoding, continuous actions and action chunks, and \pizero{}~\citep{pi0}, \pifive{}~\citep{pi05} and SmolVLA~\citep{smolvla} generate actions with a flow-matching expert. These models are compared mainly on LIBERO~\citep{libero}, simulated manipulation suites built on robosuite~\citep{robosuite} and MuJoCo~\citep{mujoco}; SimplerEnv~\citep{simplerenv} shows that simulated results can track real ones once the visual and control gaps are reduced. The evaluation code of each model is re-implemented by others, and its image steps are rarely stated in the papers.

\paragraph{Robustness evaluation.} A line of benchmarks perturbs the scene and measures the drop in success. LIBERO-Plus~\citep{liberoplus} varies \ExtPlusFactors{} factors and finds success falling from \ExtPlusFrom\% to below \ExtPlusTo\% under modest perturbations, with camera viewpoint and robot initial state the most damaging. LIBERO-PRO~\citep{liberopro} changes objects, positions, instructions and environments, and models above \ExtProFrom\% on standard LIBERO can fall to \ExtProTo\%. StableVLA~\citep{stablevla} applies \ExtStableCorruptions{} ImageNet-C corruptions~\citep{imagenetc} at severities \ExtStableSevLo{} to \ExtStableSevHi; JPEG compression at severity \ExtStableSevLo{} lowers OpenVLA on LIBERO-Spatial from \ExtStableClean\% to \ExtStableJpegThree\% in its evaluation. LIBERO-CTRL~\citep{liberoctrl} combines \ExtCtrlAxes{} perturbation axes, one of whose sensor parameters is JPEG quality, and pairs every initial state across single-axis and compound conditions; as many as \ExtCtrlFlipNull\% of matched initial states change outcome even when the net difference is not distinguishable from zero (\ExtCtrlFlipMax\% at most), so aggregate success rates hide opposing outcome changes. STRONG-VLA~\citep{strongvla} trains on \ExtStrongTypes{} perturbation types. These perturbations are large and change what the camera records; none of these studies asks whether the evaluator's own processing of a clean frame, which is meant to match training, has to be reproduced exactly. We keep the scene clean and change only those steps, whose pixel changes, apart from the geometric change of the crop, are smaller than those of the mildest ImageNet-C JPEG corruption (Appendix~\ref{app:eda}).

\paragraph{Image processing and the sensitivity of deep networks.} JPEG~\citep{wallace1992} quantizes the discrete cosine transform coefficients of $\DctBlock \times \DctBlock$ blocks, and the Lanczos filter~\citep{duchon1979} is a windowed sinc used for resampling. Image classifiers react to such operations: \citet{dodge2016} measured accuracy losses under blur, noise, contrast change and compression; ImageNet-C~\citep{imagenetc} made JPEG one of its standard corruptions; and convolutional networks change their output under small translations and rescalings because their downsampling ignores the sampling theorem~\citep{azulay2019}, which low-pass filtering before downsampling reduces~\citep{zhang2019}. Evaluation pipelines inherit this sensitivity: the resize filter and the JPEG compression inside implementations of the Fr\'echet Inception Distance change the reported scores~\citep{cleanfid}, and the JPEG quality of the images reorders model rankings on visually prompted vision-language benchmarks~\citep{vpbench}. For robot policies, compression has been studied as a transmission problem~\citep{embodiedcomp,sparc}. These studies measure open-loop outputs or heavy compression, not the closed-loop success of a policy under the mild steps of an evaluation pipeline.

\paragraph{Reproducibility and statistics of policy evaluation.} vla-eval~\citep{vlaeval} re-implements the evaluation of many VLAs in one harness; its paper reports that the undocumented crop is worth about \ExtVeCropPts{} points, and its reproduction notes list three further steps (JPEG round trip, Lanczos resize and a fixed environment seed) that were applied together with the crop and move the OpenVLA LIBERO average from \ExtVeNotesNone\% to \ExtVeNotesAll\%~\citep{vlaevalrepo}. \citet{jiang2026bench} audited LIBERO results and showed that changing only the CPU, with the GPU type fixed, changes the simulated trajectories of OpenVLA-OFT. RoboScope Kit~\citep{roboscope} compares harnesses episode by episode and finds that two runs of vla-eval itself agree on only \ExtRoboscopeSelfLo\% to \ExtRoboscopeSelfHi\% of episodes on one suite. \citet{li2026paraphrase} pair OpenVLA and OpenVLA-OFT episodes across instruction paraphrases with Wilson intervals, paired bootstrap, exact McNemar tests and Holm correction, the statistics we also use, and find that two runs of one OpenVLA-OFT configuration differ on \ExtLiRerunPct\% of episodes. A reported issue of the official OpenVLA evaluation shows that LIBERO places some fixtures differently depending on how many episodes an environment has run, which breaks pairing unless the environment loop is matched~\citep{issue342}. Concurrent work by other authors argues that success rates are not model properties and shows that the tested initial states, the instruction text and the execution horizon move VLA success by model-dependent amounts~\citep{anon2026successrates}; it does not vary the image steps of the evaluator. \citet{kressgazit2024} set out practices for policy evaluation, and \citet{agarwal2021} showed for reinforcement learning that point estimates from few runs mislead.
